% ============================================================
%  Milestone_3.tex  --  IEEE conference format (IEEEtran)
%  Not compiled by the author of this file: compile it yourself (Overleaf / pdflatex).
%
%  All measured numbers, figures and tables are generated by Milestone_3.ipynb:
%     results_macros.tex, fig_*.pdf, tab_*.tex
%  Put those next to this file. If they are missing the document still compiles and
%  shows "??" / boxes in their place, so nothing in this report is hand-typed data.
% ============================================================
\documentclass[conference]{IEEEtran}
\usepackage{cite}
\usepackage{amsmath,amssymb,amsfonts}
\usepackage{graphicx}
\usepackage{textcomp}
\usepackage{xcolor}
\usepackage{url}
\usepackage{listings}
\usepackage{microtype}
\def\BibTeX{{\rm B\kern-.05em{\sc i\kern-.025em b}\kern-.08em
    T\kern-.1667em\lower.7ex\hbox{E}\kern-.125em X}}

% ---------- safe includes (document compiles even before the notebook has been run) ----------
\newcommand{\safefig}[2]{\IfFileExists{#1}{\includegraphics[width=#2]{#1}}{\fbox{\parbox{0.85\columnwidth}{\centering\footnotesize missing #1 --- run Milestone\_3.ipynb}}}}
\newcommand{\safetab}[1]{\IfFileExists{#1}{\input{#1}}{\fbox{\footnotesize missing #1 --- run Milestone\_3.ipynb}}}

% ---------- defaults for every generated macro, then the generated values ----------
\providecommand{\bestBlocks}{\textbf{??}}
\providecommand{\bestRegs}{\textbf{??}}
\providecommand{\bwCopy}{\textbf{??}}
\providecommand{\bwNom}{\textbf{??}}
\providecommand{\ceil_regS}{\textbf{??}}
\providecommand{\ceil_regV}{\textbf{??}}
\providecommand{\ceil_smem}{\textbf{??}}
\providecommand{\coalGain}{\textbf{??}}
\providecommand{\cublasG}{\textbf{??}}
\providecommand{\cublasPctFMA}{\textbf{??}}
\providecommand{\cublasTCG}{\textbf{??}}
\providecommand{\devName}{\textbf{??}}
\providecommand{\eff_regS}{\textbf{??}}
\providecommand{\eff_regV}{\textbf{??}}
\providecommand{\eff_smem}{\textbf{??}}
\providecommand{\epsAccMax}{\textbf{??}}
\providecommand{\epsFPthirtytwoMax}{\textbf{??}}
\providecommand{\epsPosMax}{\textbf{??}}
\providecommand{\epsPosMin}{\textbf{??}}
\providecommand{\epsSignedMax}{\textbf{??}}
\providecommand{\epsSignedMin}{\textbf{??}}
\providecommand{\fmaPeak}{\textbf{??}}
\providecommand{\fmaPeakPctNom}{\textbf{??}}
\providecommand{\fuseCublasMs}{\textbf{??}}
\providecommand{\fuseCublasOnlyMs}{\textbf{??}}
\providecommand{\fuseEps}{\textbf{??}}
\providecommand{\fuseExtraMeas}{\textbf{??}}
\providecommand{\fuseExtraModel}{\textbf{??}}
\providecommand{\fuseFusedMs}{\textbf{??}}
\providecommand{\fuseN}{\textbf{??}}
\providecommand{\fuseSpeedup}{\textbf{??}}
\providecommand{\fuseVsCustom}{\textbf{??}}
\providecommand{\fuseWinN}{\textbf{??}}
\providecommand{\impliedBW}{\textbf{??}}
\providecommand{\impliedRatio}{\textbf{??}}
\providecommand{\m}{\textbf{??}}
\providecommand{\m}{\textbf{??}}
\providecommand{\m}{\textbf{??}}
\providecommand{\m}{\textbf{??}}
\providecommand{\m}{\textbf{??}}
\providecommand{\m}{\textbf{??}}
\providecommand{\m}{\textbf{??}}
\providecommand{\modelPruneSafe}{\textbf{??}}
\providecommand{\modelPrunedFrac}{\textbf{??}}
\providecommand{\modelRho}{\textbf{??}}
\providecommand{\modelTopHit}{\textbf{??}}
\providecommand{\nRef}{\textbf{??}}
\providecommand{\naiveG}{\textbf{??}}
\providecommand{\naiveOverHyp}{\textbf{??}}
\providecommand{\naivePctFMA}{\textbf{??}}
\providecommand{\nomFP}{\textbf{??}}
\providecommand{\obsClock}{\textbf{??}}
\providecommand{\pipeBestMs}{\textbf{??}}
\providecommand{\pipeBestP}{\textbf{??}}
\providecommand{\pipeBestS}{\textbf{??}}
\providecommand{\pipeDominant}{\textbf{??}}
\providecommand{\pipeHideFrac}{\textbf{??}}
\providecommand{\pipeModelErr}{\textbf{??}}
\providecommand{\pipeModelMs}{\textbf{??}}
\providecommand{\pipeN}{\textbf{??}}
\providecommand{\pipePageMs}{\textbf{??}}
\providecommand{\pipeSerialMs}{\textbf{??}}
\providecommand{\pipeSpeedup}{\textbf{??}}
\providecommand{\pipeSpeedupPage}{\textbf{??}}
\providecommand{\posDecay}{\textbf{??}}
\providecommand{\posKratio}{\textbf{??}}
\providecommand{\quantOverAcc}{\textbf{??}}
\providecommand{\regGain}{\textbf{??}}
\providecommand{\regScalarG}{\textbf{??}}
\providecommand{\regVecG}{\textbf{??}}
\providecommand{\regVecPctCublas}{\textbf{??}}
\providecommand{\regVecPctFMA}{\textbf{??}}
\providecommand{\smCount}{\textbf{??}}
\providecommand{\smemBW}{\textbf{??}}
\providecommand{\smemG}{\textbf{??}}
\providecommand{\smemGain}{\textbf{??}}
\providecommand{\tD}{\textbf{??}}
\providecommand{\tH}{\textbf{??}}
\providecommand{\tH}{\textbf{??}}
\providecommand{\tKern}{\textbf{??}}
\providecommand{\tcOverFP}{\textbf{??}}
\providecommand{\tuneBest}{\textbf{??}}
\providecommand{\tuneBestG}{\textbf{??}}
\providecommand{\tuneCount}{\textbf{??}}
\providecommand{\tuneDefaultG}{\textbf{??}}
\providecommand{\tuneGain}{\textbf{??}}
\providecommand{\tuneN}{\textbf{??}}
\providecommand{\tuneSpread}{\textbf{??}}
\providecommand{\tuneWorst}{\textbf{??}}
\providecommand{\tuneWorstG}{\textbf{??}}
\providecommand{\uncoalG}{\textbf{??}}
\providecommand{\valFail}{\textbf{??}}
\providecommand{\valRows}{\textbf{??}}
\providecommand{\valShapes}{\textbf{??}}
\providecommand{\valWorstFP}{\textbf{??}}
\providecommand{\valWorstQuant}{\textbf{??}}
\providecommand{\vecGain}{\textbf{??}}
\providecommand{\wmmaG}{\textbf{??}}
\providecommand{\wmmaPctTC}{\textbf{??}}
\IfFileExists{results_macros.tex}{\input{results_macros.tex}}{}

\lstset{basicstyle=\ttfamily\scriptsize,breaklines=true,frame=single,columns=fullflexible,xleftmargin=0pt}

\begin{document}

\title{Model-Guided Optimization of a CUDA GEMM on a Turing GPU: Where Does the Roofline Prediction Break, and When Is a Custom Kernel Worth Writing?}

\author{\IEEEauthorblockN{Shriya Soni}
\IEEEauthorblockA{2023B4AD0885G}
\and
\IEEEauthorblockN{Saee Brahmankar}
\IEEEauthorblockA{2023B4AD0909G}
\and
\IEEEauthorblockN{D Prem Kowshik}
\IEEEauthorblockA{2023A7PS0479G}
\and
\IEEEauthorblockA{\textit{Department of Computer Science} \\
\textit{BITS Pilani, Goa Campus}}}

\maketitle

\begin{abstract}
Milestone~2 produced a naive one-thread-per-element CUDA GEMM (\m2NaiveG{} GFLOP/s on an NVIDIA \devName) and conjectured that it was limited by off-chip bandwidth. This milestone replaces that conjecture with a testable analytical model and extends the project in four directions. (i)~We build an optimization ladder---shared-memory tiling, 2-D register blocking, \texttt{float4} loads and WMMA Tensor Cores---and confront a three-resource model (FMA, shared memory, DRAM) with \emph{measured} device ceilings rather than datasheet values; the model shows that the naive kernel demands \impliedRatio$\times$ the measured DRAM bandwidth, so the Milestone-2 explanation was incomplete. (ii)~We sweep 16 tile configurations and test whether the model can prune the search. (iii)~We decompose the FP16 Tensor-Core error into input-quantization and accumulation components. (iv)~We study panel-pipelined streaming as a single-GPU proxy for compute/communication overlap, and a fused GEMM+bias+ReLU epilogue as a structural argument for custom kernels. The best FP32 kernel reaches \regVecPctCublas\% of cuBLAS SGEMM at $N{=}\nRef$; the WMMA kernel reaches \wmmaPctTC\% of cuBLAS's Tensor-Core GEMM; all kernels are validated on 11 shapes including rectangular and non-tile-aligned ones.
\end{abstract}

\begin{IEEEkeywords}
GEMM, CUDA, roofline model, shared-memory tiling, register blocking, Tensor Cores, mixed precision, kernel fusion, stream pipelining
\end{IEEEkeywords}

% ============================================================
\section{Introduction}
GEMM optimization on GPUs is a mature topic: shared-memory tiling, register blocking and Tensor-Core use are textbook material, and libraries such as cuBLAS and CUTLASS already embody them. We therefore do \emph{not} claim a new GEMM algorithm. Milestones~1--2 asked how far a hand-written kernel can be pushed toward cuBLAS; that question is legitimate but, as feedback on Milestone~2 pointed out, answers itself in the literature. This milestone changes the question.

\subsection{Evolution of the research question}
Milestone~1 asked: \emph{how can a custom CUDA GEMM be optimized toward cuBLAS?} We now ask:

\begin{quote}
\textbf{Main question.} Can a simple, per-level traffic model, calibrated with microbenchmarks on the target device, \emph{predict} the performance of each optimization stage of a hand-written GEMM---and where it fails, what does the failure tell us about the hardware?
\end{quote}

\noindent It is decomposed into five measurable sub-questions:
\begin{itemize}
\item \textbf{RQ1} How well does a three-resource model (FMA throughput, shared-memory bandwidth, DRAM bandwidth) predict each stage on a Tesla T4, and was our Milestone-2 attribution of the naive kernel's limit to DRAM bandwidth correct?
\item \textbf{RQ2} How sensitive is the register-blocked kernel to its tile configuration on Turing, and can the model prune the search space?
\item \textbf{RQ3} For the FP16 Tensor-Core path, how does the error split into input quantization and FP32 accumulation, and how does it scale with $K$ and with the input distribution?
\item \textbf{RQ4} How much host--device transfer time can panel pipelining hide, and does the measured overlap follow the $\max(\cdot)$ model of Uncut-GEMMs~\cite{uncut}?
\item \textbf{RQ5} Is there a structural reason a custom kernel can still beat ``cuBLAS plus an extra pass'': a fused epilogue?
\end{itemize}

\subsection{Why a custom kernel at all?}
A fair objection is that cuBLAS exists. Our position, tested rather than asserted, has three parts. First, a library is a black box: only a from-scratch ladder lets each technique be \emph{attributed} its share of the speedup (Section~\ref{sec:ablation}) and lets a model be falsified against it. Second, tuned parameters are hardware-specific; Section~\ref{sec:tuning} checks whether Ampere-style defaults are also best on Turing. Third, fusion: a library GEMM cannot absorb arbitrary epilogue work behind its plain API, so every extra elementwise stage costs a full read--modify--write of $C$ (Section~\ref{sec:fusion}).

\subsection{Contributions}
\begin{enumerate}
\item A measured-ceiling analytical model (Eqs.~\ref{eq:bound}--\ref{eq:req}) and a quantitative re-examination of our Milestone-2 explanation.
\item An 8-stage kernel ladder with correctness checks on 11 shapes, compared against cuBLAS SGEMM and cuBLAS FP16 Tensor-Op GEMM.
\item A 16-configuration tile study with registers, spills and occupancy recorded, and a test of model-based pruning.
\item An error-budget decomposition of FP16 Tensor-Core GEMM that separates quantization from accumulation error.
\item A stream-pipelined panel GEMM validated against the $\max(\cdot)$ overlap model, and a fused-epilogue case study.
\end{enumerate}
\noindent\textbf{Scope note.} Only one GPU was available. The ``distributed scaling'' component is therefore a \emph{single-GPU proxy}: row-panel decomposition with copy/compute overlap, the intra-node analogue of tile-level scheduling in~\cite{uncut}. It is not a multi-GPU measurement and we do not present it as one.

% ============================================================
\section{Background and Related Work}
\textbf{Performance models.} The roofline model~\cite{roofline} bounds throughput by the lesser of peak compute and arithmetic intensity times bandwidth. Using \emph{compulsory} traffic, the naive GEMM looks compute-bound, which is clearly wrong; the useful move is to count traffic \emph{per memory level}, as in the microbenchmark-driven analyses of Volkov and Demmel~\cite{volkov} and the Turing dissection of Jia et al.~\cite{jiaturing}, and to calibrate bandwidths by measurement.

\textbf{Tensor Cores.} Ootomo and Yokota~\cite{ootomo23} show that shared-memory bandwidth limits Tensor-Core throughput; Choi and Stegailov~\cite{choi} show tile-size optima shifting once Tensor Cores are on; Anastasiadis et al.~\cite{uncut} hide communication via static scheduling and overlap. Markidis et al.~\cite{markidis} characterise Tensor-Core precision, and Higham and Mary~\cite{higham} survey mixed-precision linear algebra; Ootomo and Yokota~\cite{ootomo22} recover FP32 accuracy from Tensor Cores. Our RQ3 is a small, controlled instance of this error analysis tied to our own kernel.

\textbf{What is different here.} The cited works study one optimization at one scale. We keep a single kernel family on a single device and ask whether one calibrated model explains \emph{all} stages, and report where it does not.

% ============================================================
\section{Methodology}
\subsection{Platform and timing protocol}
Experiments run on an NVIDIA \devName{} (compute capability 7.5, \smCount{} SMs), CUDA~12.x, \texttt{nvcc -O3}. Milestone~2 timed one cold launch per size; here every measurement uses one untimed warm-up and 3--25 adaptive repetitions, and we report the \emph{median} (the minimum for microbenchmarks, where we want the achievable ceiling). This removes first-launch and clock-ramp artefacts that inflated small-$N$ variance in Milestone~2.

\subsection{Measured device ceilings}
Table~\ref{tab:micro} lists the ceilings used by the model, obtained with the microbenchmarks in \texttt{gemm\_cuda\_opt.cu}: a \texttt{float4} copy kernel (DRAM bandwidth), a kernel with eight independent FMA chains per thread (achievable FP32 peak) and, inside that kernel, the SM clock observed via \texttt{clock64}/\texttt{globaltimer}. Shared-memory bandwidth is estimated as $B_{\mathrm{smem}}=32\cdot4\cdot n_{\mathrm{SM}}\cdot f$ (128\,B/clk/SM at the observed clock $f$).
\begin{table}[t]
\caption{Nominal vs.\ measured device ceilings.}\label{tab:micro}
\centering\scriptsize\setlength{\tabcolsep}{3pt}
\safetab{tab_micro.tex}
\end{table}

\subsection{Analytical model}
Let $\beta_x$ be bytes moved at level $x$ per FLOP. Achievable throughput is bounded by
\begin{equation}
P \le \min\Big\{P_{\mathrm{fma}},\ \tfrac{B_{\mathrm{smem}}}{\beta_{\mathrm{smem}}},\ \tfrac{B_{\mathrm{dram}}}{\beta_{\mathrm{glob}}}\Big\}.\label{eq:bound}
\end{equation}
For a thread micro-tile $T_M\times T_N$ held in registers, each shared-memory slab read feeds $T_MT_N$ FMAs from $T_M+T_N$ loaded words, and for a block tile $B_M\times B_N$ each global word is reused across the tile:
\begin{align}
\beta_{\mathrm{smem}} &= \frac{2(T_M+T_N)}{T_M T_N}\ \text{B/FLOP},\label{eq:bsmem}\\
\beta_{\mathrm{glob}} &= 2\Big(\frac1{B_M}+\frac1{B_N}\Big)\ \text{B/FLOP}.\label{eq:bglob}
\end{align}
The shared-memory bandwidth is
\begin{equation}
B_{\mathrm{smem}} = 128\,\text{B/clk} \times n_{\mathrm{SM}} \times f_{\mathrm{obs}}.\label{eq:Bsmem}
\end{equation}
The one-element-per-thread shared-memory kernel is the special case $T_M{=}T_N{=}1$ ($\beta_{\mathrm{smem}}=4$). For the naive kernel each FMA requests two 4-byte operands, so $\beta_{\mathrm{req}}=4$\,B/FLOP at the load/store level, and a kernel running at $P$ GFLOP/s \emph{demands}
\begin{equation}
B_{\mathrm{req}} = \beta_{\mathrm{req}}\,P \quad\text{GB/s of operand traffic.}\label{eq:req}
\end{equation}
Comparing $B_{\mathrm{req}}$ with the measured $B_{\mathrm{dram}}$ tests the Milestone-2 hypothesis directly: if $B_{\mathrm{req}}>B_{\mathrm{dram}}$, requests cannot all be served by DRAM, so the limiter is cache/load-store throughput, not DRAM. The model counts bytes, not shared-memory wavefronts, and ignores bank conflicts, latency hiding and L1/L2 bandwidth; these are the expected sources of residual and are discussed in Section~\ref{sec:threats}.

\subsection{Kernel ladder}
Table~\ref{tab:ladder} describes the stages. Every FP32 kernel handles arbitrary $M,N,K$ (zero-padded tile loads, guarded stores); \texttt{float4} variants are enabled only when $K\bmod4=N\bmod4=0$.
\begin{table}[t]
\caption{Optimization ladder (pipeline steps 1--6 of Milestone~1).}\label{tab:ladder}
\centering\scriptsize\setlength{\tabcolsep}{2.5pt}
\begin{tabular}{lp{5.4cm}}
\hline
Stage & Technique \\ \hline
\texttt{naive\_m2} & 1 thread/element, $16{\times}16$ blocks, global memory only (Milestone~2) \\
\texttt{naive\_uncoalesced} & Same arithmetic, $\texttt{threadIdx.x}\to$ row: negative control for coalescing \\
\texttt{smem\_tile32} & $32{\times}32$ shared-memory tiles, 1 output/thread \\
\texttt{reg2d\_*\_scalar} & Block tile $B_M{\times}B_N{\times}B_K$, per-thread $T_M{\times}T_N$ register micro-tile, $A$ stored transposed in smem \\
\texttt{reg2d\_*\_vec} & As above with \texttt{float4} global loads \\
\texttt{wmma\_fp16\_tc} & $64{\times}64$ block tile, 4 warps, $2{\times}2$ WMMA $16{\times}16{\times}16$ fragments per warp, FP16 in / FP32 accumulate, padded smem \\
\texttt{cublas\_sgemm} & Library FP32 reference \\
\texttt{cublas\_tensorop\_fp16} & Library Tensor-Core reference (FP16 in, FP32 out) \\ \hline
\end{tabular}
\end{table}
The ``coalescing'' step of the Milestone-1 pipeline is realised as a \emph{control experiment}: the Milestone-2 kernel already maps \texttt{threadIdx.x} to the column, so loads of $B$ are coalesced and loads of $A$ are warp-broadcast; the uncoalesced variant measures how much that mapping was already worth.

\subsection{Correctness protocol}
Same Frobenius-norm criterion as Milestones~1--2, $\varepsilon_F=\|C-C_{\mathrm{ref}}\|_F/\|C_{\mathrm{ref}}\|_F$, with thresholds $\varepsilon_F<5{\times}10^{-4}$, max.\ abs.\ error $<5{\times}10^{-3}$ for FP32 kernels, checked against both the sequential \texttt{ikj} kernel and the long-double reference. The FP16 path cannot meet those thresholds against an exact reference: each input carries relative rounding error up to $u_{16}=2^{-11}$. We therefore use two checks: \emph{total} error against the original-input reference ($\varepsilon_F<2{\times}10^{-3}$, abs.\ $<5{\times}10^{-2}$) and \emph{accumulation-only} error against a reference computed on FP16-rounded inputs, which must satisfy the original FP32 thresholds. Shapes: 8 square sizes from the earlier milestones plus three rectangular, non-aligned shapes ($37{\times}129{\times}64$, $130{\times}70{\times}95$, $257{\times}33{\times}129$).

\subsection{Pipelined streaming and fused epilogue}
\emph{Pipelining.} $B$ stays resident on the device. $C$ is split into row panels of $\lceil M/P\rceil$ rows (rounded to a multiple of 128); each panel does H2D($A$ panel)$\to$kernel$\to$D2H($C$ panel) on one of $S$ non-blocking streams over pinned memory (the T4 has separate copy engines for each direction). With measured component times $T_A,T_B,T_K,T_C$ and enough streams, a three-stage pipeline predicts
\begin{equation}
T(P) = T_B + \frac{T_A{+}T_K{+}T_C}{P} + \frac{P{-}1}{P}\max(T_A,T_K,T_C).\label{eq:pipe}
\end{equation}
\emph{Fusion.} We compare $Y=\mathrm{ReLU}(AB+\mathbf{b})$ computed by (a) our register-blocked kernel with the bias+ReLU epilogue fused into the store, (b) cuBLAS SGEMM followed by a vectorised elementwise kernel, and (c) our unfused kernel followed by the same elementwise kernel (which controls for GEMM quality). The extra pass reads and writes $C$ once more:
\begin{equation}
\Delta T_{\mathrm{pass}} \approx \frac{8N^2}{B_{\mathrm{dram}}}.\label{eq:pass}
\end{equation}

\subsection{Precision hypothesis}
Rounding errors in $a_{ik},b_{kj}$ are zero-mean and independent, so for sign-symmetric data the error and the signal both grow as $\sqrt K$ and
\begin{equation}
\varepsilon_F^{\mathrm{signed}} = \mathcal{O}(u_{16}),\ \ \text{independent of }K;\qquad \varepsilon_F^{\mathrm{pos}} = \mathcal{O}\big(u_{16}/\sqrt K\big)\label{eq:err}
\end{equation}
for non-negative data, where the signal grows as $K$. FP32 accumulation adds $\mathcal{O}(\sqrt K\,u_{32})$, smaller by about $u_{32}/u_{16}=2^{-13}$. Section~\ref{sec:precision} tests these predictions.

% ============================================================
\section{Results and Discussion}
\subsection{Correctness}\label{sec:correct}
All \valRows{} kernel--shape checks over \valShapes{} shapes were run in the notebook; \valFail{} failed. The worst FP32-kernel error against the long-double reference is \valWorstFP32{}, and the worst accumulation-only error of the Tensor-Core kernel is \valWorstQuant{}. Both are at or below the FP32 rounding scale, i.e.\ two orders below the thresholds inherited from Milestone~1.

\subsection{RQ1: ablation ladder and model}\label{sec:ablation}
Fig.~\ref{fig:ablation} and Table~\ref{tab:ablation} show the ladder at $N{=}\nRef$. Relative to the Milestone-2 naive kernel (\naiveG{} GFLOP/s re-measured, i.e.\ \naivePctFMA\% of the measured FMA peak), shared-memory tiling changes throughput by $\smemGain\times$, 2-D register blocking adds a further $\regGain\times$, and \texttt{float4} loads add \vecGain\%, reaching \regVecG{} GFLOP/s (\regVecPctFMA\% of the measured FMA peak, \regVecPctCublas\% of cuBLAS SGEMM at \cublasG{} GFLOP/s). The WMMA kernel reaches \wmmaG{} GFLOP/s ($\tcOverFP32\times$ our best FP32 kernel) against \cublasTCG{} GFLOP/s for cuBLAS's Tensor-Op path. The uncoalesced control is $\coalGain\times$ slower than the Milestone-2 mapping, so a substantial part of the ``naive'' kernel's performance came from a mapping we had not examined.
\begin{figure}[t]\centering
\safefig{fig_ablation.pdf}{\columnwidth}
\caption{Throughput of every stage vs.\ problem size (log--log), with the measured FP32 FMA peak.}\label{fig:ablation}
\end{figure}
\begin{table}[t]
\caption{Ablation at $N=\nRef$. $^\dagger$Tensor-Core rows are relative to cuBLAS FP16 Tensor-Op GEMM, others to cuBLAS SGEMM. $\varepsilon_F$ is measured against cuBLAS SGEMM output.}\label{tab:ablation}
\centering\scriptsize\setlength{\tabcolsep}{2.5pt}
\safetab{tab_ablation.tex}
\end{table}

\textbf{Revisiting the Milestone-2 hypothesis.} Milestone~2 attributed the flattening of the naive kernel to ``off-chip bandwidth saturation''. Eq.~\eqref{eq:req} lets us test this: at \naiveG{} GFLOP/s the kernel demands $B_{\mathrm{req}}=\impliedBW$\,TB/s of operand traffic, which is \impliedRatio$\times$ the measured DRAM bandwidth of \bwCopy{} GB/s. If every request really went to DRAM, the kernel could not exceed \m2HypCeil{} GFLOP/s, yet it runs at \naiveOverHyp$\times$ that. The great majority of operand requests are therefore served by the L1/L2 caches, and the limiter at large $N$ is cache/load--store throughput and instruction issue, not DRAM streaming. Milestone~2's recommendation (add shared-memory tiling) still follows, but for a different reason than the one we gave: tiling replaces cache-served requests with explicitly managed on-chip reuse. We consider this correction a finding of this milestone; Nsight Compute counters in Milestone~4 will confirm the L1/L2 attribution directly.

\begin{figure}[t]\centering
\safefig{fig_model.pdf}{0.95\columnwidth}
\caption{Model ceiling (Eq.~\ref{eq:bound}) vs.\ measured throughput for the shared-memory and register-blocked kernels.}\label{fig:model}
\end{figure}
Fig.~\ref{fig:model} compares measurement with the model. The model ceilings are \ceil_smem{}, \ceil_regS{} and \ceil_regV{} TFLOP/s for the shared-memory, scalar register-blocked and vector register-blocked kernels, and the kernels achieve \eff_smem\%, \eff_regS\% and \eff_regV\% of them. The model therefore identifies the \emph{binding resource} at each stage (shared-memory bandwidth for $T_M{=}T_N{=}1$, FMA throughput once the micro-tile is $8{\times}8$), while the remaining gap to the ceiling is the unmodelled part: bank conflicts in the transposed $A$ stores, latency hiding, and the \texttt{float4}-vs-scalar load path.

\subsection{RQ2: tile-configuration sensitivity}\label{sec:tuning}
Across the \tuneCount{} configurations at $N{=}\tuneN$, throughput spans a factor \tuneSpread{} between the best (\tuneBest, \tuneBestG{} GFLOP/s, \bestRegs{} registers/thread, \bestBlocks{} blocks/SM) and the worst (\tuneWorst, \tuneWorstG{} GFLOP/s). The best configuration is \tuneGain\% faster than the Ampere-style $128{\times}128{\times}8$, $8{\times}8$ default (\tuneDefaultG{} GFLOP/s). Table~\ref{tab:tuning} lists the five best and three worst configurations with register, spill and occupancy data; Fig.~\ref{fig:tuning} overlays the model ceilings.
\begin{figure}[t]\centering
\safefig{fig_tuning.pdf}{0.95\columnwidth}
\caption{All 16 configurations ranked by measured throughput; crosses mark the model ceiling.}\label{fig:tuning}
\end{figure}
\begin{table}[t]
\caption{Best five and worst three configurations (regs = registers/thread; Ceiling = Eq.~\ref{eq:bound}, GFLOP/s).}\label{tab:tuning}
\centering\scriptsize\setlength{\tabcolsep}{2.5pt}
\safetab{tab_tuning.tex}
\end{table}

\textbf{Can the model prune the search?} The Spearman rank correlation between model ceiling and measured throughput is $\rho=\modelRho$. Configurations whose ceiling is below half the best ceiling (\modelPrunedFrac\% of the space) would be discarded without running; discarding them is ``safe'' (none within 10\% of the best measured) : \modelPruneSafe. Whether the best measured configuration lies within 10\% of the top model ceiling: \modelTopHit. Because Eq.~\eqref{eq:bound} is blind to occupancy and spills, it can eliminate poor tilings but cannot rank the good ones; that is the regime where measurement (or a richer model) is still needed.

\subsection{RQ3: error budget of the FP16 Tensor-Core path}\label{sec:precision}
Fig.~\ref{fig:precision} and Table~\ref{tab:precision} separate total error (against the original-input reference) from accumulation-only error (against the FP16-rounded-input reference). For sign-symmetric inputs the total error stays between \epsSignedMin{} and \epsSignedMax{} over $K\in[16,8192]$, on the order of $u_{16}=2^{-11}\approx4.9{\times}10^{-4}$ as predicted by Eq.~\eqref{eq:err}, while the accumulation-only error never exceeds \epsAccMax{} (FP32 kernel: \epsFPthirtytwoMax{}). Quantization therefore dominates accumulation by a median factor of \quantOverAcc{}. For non-negative inputs the total error falls from \epsPosMax{} to \epsPosMin{}, a factor \posDecay{} over a $\posKratio\times$ increase in $K$ (the $1/\sqrt K$ law predicts a factor $\approx\sqrt{\posKratio}$). \emph{Practical consequence:} on these inputs, FP16 inputs with FP32 accumulation cost a fixed $\mathcal{O}(u_{16})$ relative error regardless of $K$; improving accumulation (e.g.\ split-accumulator schemes) would not help, whereas input splitting as in~\cite{ootomo22} would.
\begin{figure}[t]\centering
\safefig{fig_precision.pdf}{0.95\columnwidth}
\caption{$\varepsilon_F$ vs.\ $K$ for FP32 and FP16-TC kernels, signed vs.\ non-negative inputs.}\label{fig:precision}
\end{figure}
\begin{table}[t]
\caption{FP16-TC error: total vs.\ accumulation-only.}\label{tab:precision}
\centering\scriptsize\setlength{\tabcolsep}{3pt}
\safetab{tab_precision.tex}
\end{table}

\subsection{RQ4: panel-pipelined streaming (single-GPU proxy)}\label{sec:pipe}
At $N{=}\pipeN$ the measured component times are H2D($A$) \tH2DA{} ms, H2D($B$) \tH2DB{} ms, kernel \tKern{} ms and D2H($C$) \tD2H{} ms; the dominant phase is the \pipeDominant. The Milestone-2-style flow (pageable memory, serial) takes \pipePageMs{} ms end-to-end; pinned memory with $P{=}1$ takes \pipeSerialMs{} ms; the best pipelined configuration ($P{=}\pipeBestP$, $S{=}\pipeBestS$) takes \pipeBestMs{} ms, a $\pipeSpeedup\times$ gain over pinned-serial and $\pipeSpeedupPage\times$ over pageable-serial, hiding \pipeHideFrac\% of the non-kernel work. Eq.~\eqref{eq:pipe} predicts \pipeModelMs{} ms at that $P$; the measurement is \pipeModelErr\% above (positive) or below (negative) the model. Fig.~\ref{fig:pipeline} and Table~\ref{tab:pipeline} show the sweep. This reproduces, at panel scale on one GPU, the effect that~\cite{uncut} exploits across GPUs: total time approaches the largest stage rather than the sum.
\begin{figure}[t]\centering
\safefig{fig_pipeline.pdf}{0.95\columnwidth}
\caption{End-to-end time vs.\ number of panels for $S=1,2,4$ streams, with the $\max(\cdot)$ model.}\label{fig:pipeline}
\end{figure}
\begin{table}[t]
\caption{End-to-end time including PCIe transfers; best $(P,S)$ per $N$.}\label{tab:pipeline}
\centering\scriptsize\setlength{\tabcolsep}{2.5pt}
\safetab{tab_pipeline.tex}
\end{table}

\subsection{RQ5: fused epilogue}\label{sec:fusion}
For $Y=\mathrm{ReLU}(AB+\mathbf b)$ at $N{=}\fuseN$ the fused custom kernel takes \fuseFusedMs{} ms, cuBLAS SGEMM plus a separate elementwise pass takes \fuseCublasMs{} ms (SGEMM alone: \fuseCublasOnlyMs{} ms), a speedup of \fuseSpeedup$\times$ of the fused kernel over that flow. Against the \emph{same} custom GEMM followed by the separate pass, fusion gives \fuseVsCustom$\times$. The extra pass costs \fuseExtraMeas{} ms measured versus \fuseExtraModel{} ms predicted by Eq.~\eqref{eq:pass}. The fused kernel is ahead of cuBLAS-plus-pass at $N\in\{\fuseWinN\}$; the fused output matches the cuBLAS-plus-pass output to $\varepsilon_F=\fuseEps$. Table~\ref{tab:fusion} and Fig.~\ref{fig:fusion} give all sizes. The structural saving ($\Delta T_{\mathrm{pass}}$) is real and predicted by the model, but it is only a net win when the custom GEMM is close enough to cuBLAS that the saved pass outweighs the GEMM deficit; the fused-versus-unfused-custom ratio isolates the saving from GEMM quality.
\begin{figure}[t]\centering
\safefig{fig_fusion.pdf}{0.95\columnwidth}
\caption{Fused vs.\ unfused execution of $\mathrm{ReLU}(AB+\mathbf b)$ (log scale).}\label{fig:fusion}
\end{figure}
\begin{table}[t]
\caption{Fused epilogue vs.\ cuBLAS + separate pass.}\label{tab:fusion}
\centering\scriptsize\setlength{\tabcolsep}{2.5pt}
\safetab{tab_fusion.tex}
\end{table}

% ============================================================
\section{Comparison with Milestone~2}
Table~\ref{tab:m2} compares the Milestone-2 results with this milestone. At $N{=}\nRef$ the best FP32 kernel is \m3OverM2Reg$\times$ the Milestone-2 kernel, the WMMA kernel \m3OverM2TC$\times$, and cuBLAS SGEMM \m3OverM2Cublas$\times$. Re-measuring the \emph{unchanged} naive kernel with the new timing protocol changes its $N{=}\nRef$ throughput by \m2RemeasureDelta\% and its $N{=}256$ throughput by \m2SmallDelta\%, quantifying the Milestone-2 cold-launch effect.
\begin{table}[t]
\caption{Throughput (GFLOP/s), Milestone~2 vs.\ Milestone~3.}\label{tab:m2}
\centering\scriptsize\setlength{\tabcolsep}{2.5pt}
\safetab{tab_m2cmp.tex}
\end{table}

% ============================================================
\section{Threats to Validity and Limitations}\label{sec:threats}
\textbf{Single, shared GPU.} Results are from one T4 whose clock may vary between runs; we report the observed clock and use it in the model, but run-to-run variation is not quantified beyond the median. \textbf{No multi-GPU hardware.} Section~\ref{sec:pipe} is a single-GPU proxy for overlap, not a scaling study. \textbf{Model scope.} Eq.~\eqref{eq:bound} counts bytes, not wavefronts; it ignores bank conflicts, L1/L2 bandwidth, latency and occupancy, and its attribution of the naive kernel's limit to cache/LSU throughput is inferred from the demanded-traffic argument, not measured with hardware counters. \textbf{Fusion baseline.} cuBLASLt offers fused bias/activation epilogues; we compare against the plain SGEMM-plus-kernel flow and do not evaluate cuBLASLt, so Section~\ref{sec:fusion} should not be read as ``custom beats the library'' in general. \textbf{Precision study.} One random seed, $M{=}N{=}128$, two input distributions. \textbf{WMMA kernel.} Scalar half loads and no double-buffering; it is a correct, illustrative Tensor-Core kernel, not a tuned one.

% ============================================================
\section{Conclusion and Next Steps}
We converted the project from ``apply known optimizations'' into a model-versus-measurement study on a Turing GPU. The calibrated three-resource model identifies the binding resource at each stage, overturned our Milestone-2 DRAM-bandwidth explanation for the naive kernel, and gives a partial but safe pruning rule for tile search; the FP16 error is quantization-dominated and follows the predicted $K$-dependence for signed versus non-negative data; and panel pipelining and epilogue fusion show where custom code can beat a plain library flow. Milestone~4 will (i) validate the per-level attributions with Nsight Compute (L1/L2 hit rates, shared-memory bank conflicts, issue-slot utilisation), (ii) add double-buffered, bank-conflict-free tiling and compare against cuBLASLt epilogues, and (iii) extend the pipelined study to multiple GPUs if hardware is available.

\section*{Reproducibility}
\texttt{nvcc -O3 -arch=sm\_75 -std=c++14 -lineinfo -o gemm\_cuda\_opt gemm\_cuda\_opt.cu -lcublas}; run \texttt{./gemm\_cuda\_opt \{micro|validate|bench|tune|precision|pipe|fuse|all\}}. \texttt{Milestone\_3.ipynb} runs every mode and generates all figures, tables and the macro file used by this document.

\begin{thebibliography}{99}
\bibitem{ootomo23} H.~Ootomo and R.~Yokota, ``Reducing shared memory footprint to leverage high throughput on Tensor Cores and its flexible API extension library,'' ACM, 2023, doi:10.1145/3578178.3578238.
\bibitem{choi} Y.~R. Choi and V.~Stegailov, ``Multi-GPU GEMM algorithm performance analysis for Nvidia and AMD GPUs connected by NVLink and PCIe,'' in \emph{Mathematical Modeling and Supercomputer Technologies (MMST 2022)}, CCIS vol.~1750, Springer, 2022, pp.~281--292.
\bibitem{uncut} P.~Anastasiadis, N.~Papadopoulou, N.~Koziris, and G.~Goumas, ``Uncut-GEMMs: Communication-aware matrix multiplication on multi-GPU nodes,'' in \emph{2024 IEEE Int.\ Conf.\ Cluster Computing (CLUSTER)}, 2024, pp.~143--154.
\bibitem{roofline} S.~Williams, A.~Waterman, and D.~Patterson, ``Roofline: An insightful visual performance model for multicore architectures,'' \emph{Commun.\ ACM}, vol.~52, no.~4, pp.~65--76, 2009.
\bibitem{volkov} V.~Volkov and J.~W. Demmel, ``Benchmarking GPUs to tune dense linear algebra,'' in \emph{Proc.\ SC'08}, 2008.
\bibitem{jiaturing} Z.~Jia, M.~Maggioni, J.~Smith, and D.~P. Scarpazza, ``Dissecting the NVIDIA Turing T4 GPU via microbenchmarking,'' arXiv:1903.07486, 2019.
\bibitem{markidis} S.~Markidis, S.~W.~D. Chien, E.~Laure, I.~B. Peng, and J.~S. Vetter, ``NVIDIA Tensor Core programmability, performance \& precision,'' in \emph{IPDPS Workshops}, 2018.
\bibitem{higham} N.~J. Higham and T.~Mary, ``Mixed precision algorithms in numerical linear algebra,'' \emph{Acta Numerica}, vol.~31, pp.~347--414, 2022.
\bibitem{ootomo22} H.~Ootomo and R.~Yokota, ``Recovering single precision accuracy from Tensor Cores while surpassing the FP32 theoretical peak performance,'' \emph{Int.\ J.\ High Perform.\ Comput.\ Appl.}, 2022.
\end{thebibliography}

\end{document}
